% CSE 715 — Computer Graphics
% Chapter 6: Three-Dimensional Transformations — Complete Model Answers (2020–2024)
% University of Chittagong, 7th Semester B.Sc. Engineering
% Source: Schaum's Outline of Computer Graphics (Xiang & Plastock), Chapter 6
%
% Compile with: pdflatex Chapter6_Solutions.tex (run twice for TOC)

\documentclass[12pt, a4paper]{article}

\usepackage[left=2.5cm, right=2.5cm, top=2.8cm, bottom=2.8cm, headheight=15pt]{geometry}
\usepackage{amsmath, amssymb}
\usepackage{tikz}
\usetikzlibrary{arrows.meta, positioning, shapes.geometric, calc, backgrounds, fit}
\usepackage{booktabs}
\usepackage{array}
\usepackage{xcolor}
\usepackage{enumitem}
\usepackage{fancyhdr}
\usepackage{titlesec}
\usepackage{mdframed}
\usepackage{multicol}
\usepackage{tabularx}
\usepackage{multirow}
\usepackage{hyperref}

%----- Colours -----
\definecolor{sectionblue}{RGB}{10,60,110}
\definecolor{boxbg}{RGB}{242,247,255}
\definecolor{answerbg}{RGB}{240,255,240}
\definecolor{notesbg}{RGB}{255,250,230}
\definecolor{keyblue}{RGB}{0,80,160}
\definecolor{accent}{RGB}{180,40,40}

%----- Box environments -----
\mdfdefinestyle{solutionframe}{linecolor=sectionblue, backgroundcolor=boxbg, roundcorner=4pt, linewidth=1.4pt, innertopmargin=7pt, innerbottommargin=7pt, innerleftmargin=9pt, innerrightmargin=9pt}
\mdfdefinestyle{answerframe}{linecolor=green!55!black, backgroundcolor=answerbg, roundcorner=4pt, linewidth=1.4pt, innertopmargin=7pt, innerbottommargin=7pt, innerleftmargin=9pt, innerrightmargin=9pt}
\mdfdefinestyle{notesframe}{linecolor=orange!70!black, backgroundcolor=notesbg, roundcorner=4pt, linewidth=1pt, innertopmargin=5pt, innerbottommargin=5pt, innerleftmargin=8pt, innerrightmargin=8pt}
\newenvironment{solutionbox}{\begin{mdframed}[style=solutionframe]}{\end{mdframed}}
\newenvironment{answerbox}{\begin{mdframed}[style=answerframe]}{\end{mdframed}}
\newenvironment{notesbox}{\begin{mdframed}[style=notesframe]}{\end{mdframed}}

%----- Page style -----
\pagestyle{fancy}
\fancyhf{}
\lhead{\small\color{sectionblue} CSE 715 --- Chapter 6: 3D Transformations}
\rhead{\small\color{sectionblue} Model Solutions (2020--2024)}
\cfoot{\small\thepage}
\renewcommand{\headrulewidth}{0.5pt}

%----- Section formatting -----
\titleformat{\section}{\Large\bfseries\color{sectionblue}}{}{0em}{}[\vspace{-2pt}\color{sectionblue}\rule{\linewidth}{0.8pt}]
\titleformat{\subsection}{\large\bfseries\color{sectionblue!85!black}}{}{0em}{}
\titleformat{\subsubsection}{\normalsize\bfseries\color{black}}{}{0em}{}

%----- Custom list -----
\newlist{keypoints}{itemize}{1}
\setlist[keypoints]{label=$\triangleright$, leftmargin=1.5cm, itemsep=2pt, topsep=3pt}
\newlist{examlist}{enumerate}{1}
\setlist[examlist]{label=\textbf{(\roman*)}, leftmargin=1.8cm, itemsep=3pt, topsep=4pt}

%=====================================================================
\begin{document}
%=====================================================================

\begin{center}
  {\LARGE\bfseries\color{sectionblue} CSE 715: Computer Graphics}\\[0.6em]
  {\Large\bfseries Chapter 6 --- Three-Dimensional Transformations}\\[0.3em]
  {\large Complete Model Answers (2020--2024)}\\[0.2em]
  {\normalsize University of Chittagong \quad|\quad 7\textsuperscript{th} Semester B.Sc.\ (Engg.) Examination}\\[0.3em]
  {\small Source: Schaum's Outline of Computer Graphics (Xiang \& Plastock), Chapter 6 --- notation follows the book exactly ($R_{\theta,\mathbf{I}}$/$R_{\theta,\mathbf{J}}$/$R_{\theta,\mathbf{K}}$ for rotation about $x/y/z$, $S_{s_x,s_y,s_z}$ for scaling, $T_{\mathbf{V}}$ for translation)}
\end{center}

\vspace{0.3em}
\noindent\rule{\linewidth}{1pt}
\vspace{0.5em}

\tableofcontents
\newpage

%=====================================================================
\section{Theoretical Framework --- Reusable Across All Years}
%=====================================================================
\begin{notesbox}
\textbf{Yield note:} Chapter 6 is the lightest exam chapter --- across all 5 years (2020--2024), only \textbf{one} genuine Chapter-6 question appears: \textbf{tilting} (2022, Q5c, 3.5 marks). \textbf{3D scaling} (2023, Q3c, 4 marks) is Chapter-4 scaling machinery applied to 3D points and is already solved in Chapter4\_Solutions.pdf --- it is re-solved here too since it is Ch6 material by chapter number. \textbf{Mirror reflection} is syllabus-mandated ("Chapter 6: rotation, transformation, scaling, mirror") but has \textbf{0/5 PYQ hits} --- included below for completeness only, same treatment as Bezier curves in Ch9.
\end{notesbox}

%----------------------------------------------------------------------
\subsection{Translation, Scaling, and Canonical-Axis Rotation (\S6.1)}
%----------------------------------------------------------------------
\begin{answerbox}
\textbf{Translation} by $\mathbf{V}=a\mathbf{I}+b\mathbf{J}+c\mathbf{K}$: $x'=x+a,\ y'=y+b,\ z'=z+c$.
$$T_{\mathbf{V}} = \begin{pmatrix}1&0&0&a\\0&1&0&b\\0&0&1&c\\0&0&0&1\end{pmatrix}$$
\textbf{Scaling about the origin} with factors $s_x,s_y,s_z$ ($s>1$ magnifies, $s<1$ reduces): $x'=s_x\cdot x,\ y'=s_y\cdot y,\ z'=s_z\cdot z$.
$$S_{s_x,s_y,s_z} = \begin{pmatrix}s_x&0&0\\0&s_y&0\\0&0&s_z\end{pmatrix}$$
\textbf{Canonical rotations} (right-hand rule, positive angle $\theta$ about the named axis):
$$R_{\theta,\mathbf{K}}\ (\text{about }z):\begin{cases}x'=x\cos\theta-y\sin\theta\\y'=x\sin\theta+y\cos\theta\\z'=z\end{cases} \qquad R_{\theta,\mathbf{J}}\ (\text{about }y):\begin{cases}x'=x\cos\theta+z\sin\theta\\y'=y\\z'=-x\sin\theta+z\cos\theta\end{cases} \qquad R_{\theta,\mathbf{I}}\ (\text{about }x):\begin{cases}x'=x\\y'=y\cos\theta-z\sin\theta\\z'=y\sin\theta+z\cos\theta\end{cases}$$
$$R_{\theta,\mathbf{K}}=\begin{pmatrix}\cos\theta&-\sin\theta&0\\\sin\theta&\cos\theta&0\\0&0&1\end{pmatrix} \quad R_{\theta,\mathbf{J}}=\begin{pmatrix}\cos\theta&0&\sin\theta\\0&1&0\\-\sin\theta&0&\cos\theta\end{pmatrix} \quad R_{\theta,\mathbf{I}}=\begin{pmatrix}1&0&0\\0&\cos\theta&-\sin\theta\\0&\sin\theta&\cos\theta\end{pmatrix}$$
\textbf{Coordinate transformations} (move the observer, object fixed) use the same matrices with $\theta\to-\theta$ and $\mathbf{V}\to-\mathbf{V}$ (e.g.\ $\bar T_{\mathbf V}=T_{-\mathbf V}$, $\bar R_\theta=R_{-\theta}$) --- direct 3D analogue of Ch4's geometric-vs-coordinate distinction. \textbf{Homogeneous form:} all 3$\times$3 matrices above adjoin a 4th row/column ($0\,0\,0\,1$) so rotation, scaling, and translation can be concatenated by matrix multiplication (order matters --- rightmost matrix applies first to a column vector).
\end{answerbox}

%----------------------------------------------------------------------
\subsection{Composite / Tilting Transformation (\S6.3, Solved Problem 6.1)}
%----------------------------------------------------------------------
\begin{answerbox}
\textbf{Tilting} is defined as a rotation about the $x$-axis (angle $\theta_x$) followed by a rotation about the $y$-axis (angle $\theta_y$). Since the $x$-rotation is applied first, it sits closest to the vector being transformed:
$$T = R_{\theta_y,\mathbf{J}}\cdot R_{\theta_x,\mathbf{I}} = \begin{pmatrix}\cos\theta_y & \sin\theta_y\sin\theta_x & \sin\theta_y\cos\theta_x\\ 0 & \cos\theta_x & -\sin\theta_x\\ -\sin\theta_y & \cos\theta_y\sin\theta_x & \cos\theta_y\cos\theta_x\end{pmatrix}$$
\textbf{Does order matter?} Reversing the order (rotate about $y$ first, then $x$): $R_{\theta_x,\mathbf{I}}\cdot R_{\theta_y,\mathbf{J}} = \begin{pmatrix}\cos\theta_y & 0 & \sin\theta_y\\ \sin\theta_x\sin\theta_y & \cos\theta_x & -\sin\theta_x\cos\theta_y\\ -\cos\theta_x\sin\theta_y & \sin\theta_x & \cos\theta_x\cos\theta_y\end{pmatrix}$ --- \textbf{not the same matrix as $T$} (compare position (2,1): $0$ vs.\ $\sin\theta_x\sin\theta_y$). \textbf{Matrix multiplication does not commute in general $\Rightarrow$ order matters} for 3D composite rotations, unlike scalar multiplication.
\end{answerbox}

%----------------------------------------------------------------------
\subsection{Rotation About an Arbitrary Axis --- Alignment (\S6.3, Solved Problems 6.2--6.4)}
%----------------------------------------------------------------------
\begin{answerbox}
To rotate by $\theta°$ about an arbitrary axis $L$ (direction $\mathbf V=a\mathbf I+b\mathbf J+c\mathbf K$, through point $P$): (1) translate $P$ to the origin ($T_{-P}$); (2) align $\mathbf V$ with $\mathbf K$ via $A_{\mathbf V}$ (two canonical rotations: about $x$ by $\theta_1$ with $\sin\theta_1=\frac{b}{\sqrt{b^2+c^2}},\cos\theta_1=\frac{c}{\sqrt{b^2+c^2}}$, then about $y$ by $-\theta_2$ with $\lambda=\sqrt{b^2+c^2}$, $\sin(-\theta_2)=\frac{-a}{|\mathbf V|},\cos(-\theta_2)=\frac{\lambda}{|\mathbf V|}$); (3) rotate by $\theta°$ about $\mathbf K$ ($R_{\theta,\mathbf K}$); (4) reverse steps 2 and 1:
$$R_{\theta,L} = T_{-P}^{-1}\cdot A_{\mathbf V}^{-1}\cdot R_{\theta,\mathbf K}\cdot A_{\mathbf V}\cdot T_{-P}$$
\end{answerbox}

%----------------------------------------------------------------------
\subsection{Mirror Reflection (\S6.3, Solved Problems 6.6--6.8) --- syllabus-mandated, 0/5 PYQ}
%----------------------------------------------------------------------
\begin{answerbox}
\textbf{Reflection about the $xy$-plane:} $P(x,y,z)\to P'(x,y,-z)$, i.e.\ $M=\begin{pmatrix}1&0&0\\0&1&0\\0&0&-1\end{pmatrix}$.
\textbf{Reflection about an arbitrary plane} through $P_0(x_0,y_0,z_0)$ with normal $\mathbf N$: translate $P_0$ to the origin, align $\mathbf N$ with $\mathbf K$, apply the $xy$-plane reflection $M$, then reverse: $M_{\mathbf N,P_0}=T_{\mathbf V}^{-1}\cdot A_{\mathbf N}^{-1}\cdot M\cdot A_{\mathbf N}\cdot T_{\mathbf V}$ where $\mathbf V=-x_0\mathbf I-y_0\mathbf J-z_0\mathbf K$.
\end{answerbox}

%=====================================================================
\section{Examination 2022 --- Q5(c) (Section B): Tilting}
%=====================================================================
\begin{solutionbox}
\textbf{Question:} Define tilting as a rotation about the $x$-axis followed by a rotation about the $y$-axis. (i) Find the tilting matrix. (ii) Does the order of performing the rotation matter? [3.5 marks]
\end{solutionbox}

\subsection*{Solution}
\begin{solutionbox}
\textbf{Definition:} tilting = rotate about the $x$-axis by $\theta_x$, then rotate the result about the $y$-axis by $\theta_y$.

\textbf{(i) Tilting matrix} --- $x$-rotation applied first, so it is rightmost in the product:
$$T = R_{\theta_y,\mathbf{J}}\cdot R_{\theta_x,\mathbf{I}} = \begin{pmatrix}\cos\theta_y & 0 & \sin\theta_y\\ 0 & 1 & 0\\ -\sin\theta_y & 0 & \cos\theta_y\end{pmatrix}\begin{pmatrix}1&0&0\\0&\cos\theta_x&-\sin\theta_x\\0&\sin\theta_x&\cos\theta_x\end{pmatrix} = \begin{pmatrix}\cos\theta_y & \sin\theta_y\sin\theta_x & \sin\theta_y\cos\theta_x\\ 0 & \cos\theta_x & -\sin\theta_x\\ -\sin\theta_y & \cos\theta_y\sin\theta_x & \cos\theta_y\cos\theta_x\end{pmatrix}$$

\textbf{(ii) Does order matter?} Reverse the order --- rotate about $y$ first, then $x$:
$$R_{\theta_x,\mathbf{I}}\cdot R_{\theta_y,\mathbf{J}} = \begin{pmatrix}\cos\theta_y & 0 & \sin\theta_y\\ \sin\theta_x\sin\theta_y & \cos\theta_x & -\sin\theta_x\cos\theta_y\\ -\cos\theta_x\sin\theta_y & \sin\theta_x & \cos\theta_x\cos\theta_y\end{pmatrix}$$
This differs from $T$ (e.g.\ entry $(2,1)$ is $0$ in $T$ but $\sin\theta_x\sin\theta_y$ here). \textbf{Yes, the order matters} --- matrix multiplication is not commutative, so "rotate about $x$ then $y$" and "rotate about $y$ then $x$" give physically different final orientations.
\end{solutionbox}

%=====================================================================
\section{Examination 2023 --- Q3(c) (Section A): 3D Scaling (Related --- also in Chapter 4)}
%=====================================================================
\begin{solutionbox}
\textbf{Question:} Given a 3D object with coordinate points $A(3,0,3)$, $B(3,3,6)$, $C(3,0,1)$, and $D(0,0,0)$. Apply the scaling parameter 3 towards the $X$ axis, 2 towards the $Y$ axis, and 3 towards the $Z$ axis, and obtain the new coordinates of the object. [4 marks]
\end{solutionbox}

\subsection*{Solution}
\begin{solutionbox}
Scaling about the origin with $s_x=3,\ s_y=2,\ s_z=3$: $x'=3x,\ y'=2y,\ z'=3z$.
$$S_{3,2,3}=\begin{pmatrix}3&0&0\\0&2&0\\0&0&3\end{pmatrix}$$
\begin{center}
\begin{tabular}{lccc}
\toprule
Vertex & Original $(x,y,z)$ & Apply $(3x,2y,3z)$ & New $(x',y',z')$ \\
\midrule
$A$ & $(3,0,3)$ & $(9,0,9)$ & $\mathbf{A'(9,0,9)}$ \\
$B$ & $(3,3,6)$ & $(9,6,18)$ & $\mathbf{B'(9,6,18)}$ \\
$C$ & $(3,0,1)$ & $(9,0,3)$ & $\mathbf{C'(9,0,3)}$ \\
$D$ & $(0,0,0)$ & $(0,0,0)$ & $\mathbf{D'(0,0,0)}$ \\
\bottomrule
\end{tabular}
\end{center}
Identical result to Chapter4\_Solutions.pdf's treatment of this same question (2D-scaling machinery generalizes directly --- each axis just scales independently).
\end{solutionbox}

%=====================================================================
\section{Quick Summary --- Exam Patterns for Chapter 6}
%=====================================================================

\begin{center}
\renewcommand{\arraystretch}{1.3}
\begin{tabular}{p{4.5cm}|p{2.4cm}|p{7.1cm}}
\toprule
\textbf{Pattern} & \textbf{Years} & \textbf{Method to memorise} \\
\midrule
Tilting (composite rotation: $x$ then $y$) + does order matter & 2022 only & $T=R_{\theta_y,\mathbf J}\cdot R_{\theta_x,\mathbf I}$; multiply the two canonical matrices, plug in, then show the reversed product differs $\Rightarrow$ order matters \\
3D scaling application about the origin & 2023 only & Multiply each coordinate independently by its axis scale factor $s_x,s_y,s_z$ --- identical to 2D scaling, one more axis \\
Mirror reflection (syllabus-mandated, never asked) & 0/5 & $xy$-plane: negate $z$. Arbitrary plane: translate-align-reflect-reverse, same recipe as arbitrary-axis rotation \\
\bottomrule
\end{tabular}
\end{center}

\begin{notesbox}
\textbf{Exam-day strategy for Chapter 6:} \textbf{correction from earlier tracking} --- tilting was previously logged as a 2/5-year topic (2021+2022); re-verified directly against all 5 raw scanned papers on 2026-07-05 and it only appears in \textbf{2022 (Q5c)}. Given the entire chapter is worth at most $\sim$4 marks in a typical year and sometimes 0, this is a "learn the method once, move on" chapter --- don't over-invest revision time here relative to Ch5/Ch7/Ch10.
\end{notesbox}

\end{document}
